\documentclass[11pt,a4paper]{article}
\usepackage[utf8]{inputenc}
\usepackage[T1]{fontenc}
\usepackage[margin=2.5cm]{geometry}
\usepackage{amsmath,amssymb,amsthm}
\usepackage{booktabs,array}
\usepackage{xcolor}
\usepackage{tcolorbox}
\tcbuselibrary{skins,breakable}
\usepackage{tikz}
\usetikzlibrary{arrows.meta,positioning,shapes.geometric,calc,fit,backgrounds}
\usepackage{hyperref}
\hypersetup{colorlinks=true,linkcolor=blue!60!black,urlcolor=blue!60!black,
            citecolor=blue!60!black}
\setlength{\emergencystretch}{2em}
\usepackage{enumitem}
\usepackage[numbers,sort&compress]{natbib}

\definecolor{boxblue}{RGB}{220,235,255}
\definecolor{boxorange}{RGB}{255,240,210}
\definecolor{boxgreen}{RGB}{220,245,220}
\definecolor{boxgold}{RGB}{255,248,210}

\newcommand{\term}[1]{\textbf{#1}}
\newcommand{\lean}[1]{%
  \penalty-100
  \begingroup\def\_{\textunderscore\allowbreak}\texttt{#1}\endgroup}
\newcommand{\CatAL}{\ensuremath{\mathbf{A}_{\rm Lean}}}
\newcommand{\CatAD}{\textbf{A/D}}
\newcommand{\CatA}{\textbf{A}}

\newtheorem{theorem}{Theorem}[section]
\newtheorem{result}[theorem]{Result}

\newcommand{\TutorialDOI}{10.5281/zenodo.20669146}
\date{2026}

\title{\textbf{The Hierarchy Problem Dissolved:}\\[6pt]
\large Deriving the Higgs Mass and Electroweak Scale from the GTE Polynomial\\[8pt]
\normalsize
\href{https://novaspivack.com/research}{novaspivack.com/research}\quad
\href{https://doi.org/10.5281/zenodo.20560550}{P48 complete theory monograph}\\[4pt]
\small\href{https://doi.org/\TutorialDOI}{doi.org/\TutorialDOI}
  \textit{(registration pending)}\\[4pt]
\normalsize\textit{UGP Physics Tutorial Series}}
\author{Nova Spivack}

\begin{document}
\setcounter{tocdepth}{2}
\maketitle
\tableofcontents
\newpage

\begin{tcolorbox}[colback=blue!6,colframe=blue!40!black,
  title=\textbf{UGP Physics Tutorial Series}]
\textbf{Prerequisites (read first):}
  \href{https://doi.org/10.5281/zenodo.20657868}{\emph{The Master Quadratic: One Equation, Two Crown Jewels}}
  $\cdot$
  \href{https://doi.org/10.5281/zenodo.20657885}{\emph{The $\Phi_{\rm MDL}$ Field: From Discrete Certificate to Continuum Physics}}\\[4pt]
\textbf{Next tutorials:}
  \href{https://doi.org/10.5281/zenodo.20669144}{\emph{The Forces: How SU(3)$\times$SU(2)$\times$U(1) Emerges from the Polynomial}}\\[4pt]
\textbf{See also:}
  \href{https://doi.org/10.5281/zenodo.20657876}{\emph{How MDL Selects the 19-Bit Polynomial}}
  $\cdot$
  \href{https://doi.org/10.5281/zenodo.20657889}{\emph{The GTE Polynomial: A Step-by-Step Cheat Sheet}}
\end{tcolorbox}

\begin{tcolorbox}[colback=boxgreen,colframe=green!50!black,
  title=\textbf{Claim-Strength Taxonomy}]
Throughout this tutorial, results are labeled with a confidence grade:
\begin{itemize}[itemsep=2pt]
  \item \textbf{$\mathbf{A}_{\rm Lean}$ (CatAL)} --- Machine-certified in Lean~4 with zero \texttt{sorry} and zero custom axioms.  Independently verifiable by anyone with Lean installed.
  \item \textbf{A/D (CatAD)} --- Analytically derived with full derivation, computationally confirmed.  Not yet machine-certified.
  \item \textbf{A (CatA)} --- Analytically derived; derivation complete but not independently machine-certified.
  \item \textbf{A/D (conditional)} --- Derived under a named, stated premise; the premise is itself an open problem or a CatB assumption.
  \item \textbf{B (CatB)} --- A stated working hypothesis or structural premise; not yet derived.
\end{itemize}
\end{tcolorbox}

%─────────────────────────────────────────────────────────────────────────────
\section{The Hierarchy Problem: What It Is and Why It Matters}
\label{sec:problem}
%─────────────────────────────────────────────────────────────────────────────

\begin{tcolorbox}[colback=boxblue,colframe=blue!40!black,title=The question in one paragraph]
The Higgs boson weighs 125~GeV.
The Planck scale --- where quantum gravity becomes important --- is at
$10^{19}$~GeV.
That is 17 orders of magnitude apart.
Quantum mechanics says the Higgs mass should be pulled up to the Planck
scale by quantum corrections.
Why is the Higgs so light?
The Standard Model has no answer.
The GTE framework does: the electroweak scale is a fixed point of the
polynomial's own arithmetic.
\end{tcolorbox}

\subsection{Two scales with no bridge}

Physics has two fundamental scales.
The \term{electroweak scale} sets the masses of the $W$ and $Z$ bosons
($\sim 80$--$91$~GeV) and the Higgs boson ($m_H = 125.20 \pm 0.11$~GeV,
PDG~2024~\cite{SpivackGTECompleteFramework}).
It is also the scale of the \term{Higgs vacuum expectation value} (VEV)
$v = 246.22$~GeV, which is the field-theory quantity from which the
$W$, $Z$, and Higgs masses are derived.

The \term{Planck scale} $M_{\rm Pl} = 1.22 \times 10^{19}$~GeV is where
the gravitational and quantum effects are expected to become comparable.
The ratio of these scales is:
\[
  \frac{M_{\rm Pl}}{v} \approx 5 \times 10^{16}.
\]
Seventeen orders of magnitude separate the two.

In the Standard Model (SM), the Higgs VEV $v = 246.22$~GeV is a free parameter
inserted by hand.
The SM says nothing about why this particular value, rather than
$10^{19}$~GeV or $10^{-35}$~GeV, is the physical one.

\subsection{Why quantum corrections make this worse}

In quantum field theory, every particle mass receives \term{quantum corrections}:
contributions from virtual particles that loop through Feynman diagrams.
For fermion masses (quarks and leptons), these corrections are proportional
to the mass itself, so they are naturally small --- a phenomenon called
\term{chiral symmetry protection}.

The Higgs boson is a scalar (spin-0) particle.
Scalars do not benefit from chiral protection.
Their quantum corrections are proportional to the highest energy scale in
the theory, not to the mass itself.
Concretely, the one-loop correction to the Higgs mass-squared is:
\[
  \delta m_H^2 \sim \frac{\Lambda_{\rm UV}^2}{16\pi^2},
\]
where $\Lambda_{\rm UV}$ is the ultraviolet cutoff --- the highest energy
at which the SM is valid.
If the SM is valid up to the Planck scale, then
$\delta m_H^2 \sim (10^{19}~{\rm GeV})^2 / (16\pi^2)$,
which is $(10^{18}~{\rm GeV})^2$ --- larger than the observed Higgs mass
by 34 orders of magnitude.

To get $m_H = 125$~GeV from an underlying theory at the Planck scale,
the SM must cancel this enormous correction against an equally enormous
bare Higgs mass parameter, with 32 decimal places of coincidence.
This is called \term{fine-tuning}, and it is considered the
deepest naturalness problem in particle physics.

\subsection{A worked example of the fine-tuning}

Let $m_{H,0}$ be the ``bare'' Higgs mass parameter in the SM Lagrangian.
The physical Higgs mass is:
\[
  m_H^2 = m_{H,0}^2 + \delta m_H^2.
\]
We want $m_H^2 = (125~{\rm GeV})^2 = 1.56 \times 10^4~{\rm GeV}^2$.
The quantum correction (for $\Lambda_{\rm UV} = M_{\rm Pl}$) is
$\delta m_H^2 \approx -(3/8\pi^2) y_t^2 M_{\rm Pl}^2 \approx -5 \times 10^{36}~{\rm GeV}^2$
(dominated by the top quark loop).

Therefore $m_{H,0}^2$ must equal $+5 \times 10^{36}~{\rm GeV}^2$ to 32
decimal places.
Two enormous numbers, each of order $10^{36}$, cancel to give $10^4$.
This is an extraordinary coincidence --- one that demands explanation.

%─────────────────────────────────────────────────────────────────────────────
\section{Standard Solutions and Why They Fall Short}
\label{sec:standard}
%─────────────────────────────────────────────────────────────────────────────

Five standard approaches to the hierarchy problem have been proposed
and tested.

\subsection{Supersymmetry}

\term{Supersymmetry} (SUSY) introduces a \term{superpartner} for every SM particle:
a boson for each fermion and vice versa.
The quantum corrections from a boson and its superpartner have opposite sign
and cancel exactly, eliminating the fine-tuning problem.

The problem: supersymmetric particles have not been observed.
The Large Hadron Collider (LHC) has searched for sparticles up to
$\sim 2$~TeV and found none.
The minimal SUSY models are under severe experimental pressure.
The more parameters are added to push sparticles heavier, the more SUSY
itself becomes a fine-tuned theory.

\subsection{Extra dimensions}

Models with extra spatial dimensions (Randall--Sundrum, ADD) can make
the Planck scale appear large from our perspective while the ``true''
Planck scale in the full higher-dimensional spacetime is close to the
electroweak scale.
No experimental evidence for extra dimensions has been found.

\subsection{Technicolor}

\term{Technicolor} models replaced the Higgs scalar with a condensate of
new strongly-interacting fermions, analogous to how quark condensates
generate the pion mass.
Precision electroweak measurements at LEP ruled out minimal Technicolor
models in the 1990s.

\subsection{Relaxion}

The \term{relaxion} is a new scalar field whose value rolls slowly during
cosmological evolution, dynamically selecting the electroweak scale.
No observational evidence supports the relaxion.
The mechanism introduces new arbitrary parameters.

\subsection{Landscape/anthropic}

The string landscape contains $10^{500}$ vacuum states.
Some have Higgs masses near 125~GeV.
The anthropic argument asserts that only those vacua support observers,
so we necessarily observe a light Higgs.
This is unfalsifiable by construction.

\medskip

The GTE framework takes a structurally different approach.
It does not cancel, relax, or anthropically select the Higgs mass.
It derives it as the fixed point of the polynomial's own equation.

%─────────────────────────────────────────────────────────────────────────────
\section{The GTE Approach: The SRRG Fixed Point}
\label{sec:srrg}
%─────────────────────────────────────────────────────────────────────────────

\begin{tcolorbox}[colback=boxgreen,colframe=green!50!black,
  title=The key idea]
The GTE polynomial $p(L,C,R)$ contains its own self-referential
renormalization group equation.
The fixed point of that equation --- the coupling value that maps to itself ---
is the golden ratio $\varphi = (\sqrt{5}-1)/2$.
The electroweak scale is set by this fixed point.
No cancellation. No new particles. No fine-tuning.
\end{tcolorbox}

\subsection{The master quadratic hidden in the polynomial}

The GTE polynomial is $p(L,C,R) = C + R - CR - LCR \pmod{7}$
(all arithmetic mod 7; see \emph{The GTE Polynomial: A Step-by-Step
Cheat Sheet}~\cite{SpivackGTECompleteFramework,SpivackTutorialPolynomial}).

A natural question: what happens when all three inputs are the same?
Setting $L = C = R = x$ (``evaluating on the diagonal''):
\[
  p(x,x,x) = x + x - x^2 - x^3 = 2x - x^2 - x^3.
\]
This is no longer a function of three separate inputs; it is a polynomial in $x$ alone.

Now ask: for which $x$ does a uniform state remain uniform?
A uniform cellular automaton configuration (every cell equal to $x$)
stays uniform when $p(x,x,x) = x$.
Setting the diagonal output equal to the input:
\[
  2x - x^2 - x^3 = x \implies x - x^2 - x^3 = 0
  \implies -x\bigl(x^2 + x - 1\bigr) = 0.
\]
There are two solutions:
\begin{enumerate}[itemsep=2pt]
  \item $x = 0$ (the trivial vacuum, every cell zero).
  \item $x^2 + x - 1 = 0$ (a non-trivial equation).
\end{enumerate}

\subsection{The non-trivial fixed point: the golden ratio}

The equation $x^2 + x - 1 = 0$ has two solutions:
\[
  x = \frac{-1 \pm \sqrt{5}}{2}.
\]
The unique positive real solution is:
\[
  x^* = \frac{\sqrt{5}-1}{2} = \varphi \approx 0.61803\ldots,
\]
which is the \term{golden ratio} $\varphi$.
The unique negative real solution is $-(1+\varphi) = -1/\varphi$.

The golden ratio satisfies $\varphi^2 + \varphi = 1$,
which is exactly the fixed-point equation written differently.
This is not a coincidence inserted by hand --- it is forced by the polynomial's
own structure.

\subsection{The Self-Referential Renormalization Group}

The \term{Self-Referential Renormalization Group} (SRRG) is the renormalization
group whose $\beta$-function equation is the GTE polynomial applied to
its own coupling~\cite{SpivackSRRG}:
\[
  \beta(g) = 0 \iff g = p(g,g,g).
\]
A coupling $g^*$ is a fixed point if and only if $p(g^*,g^*,g^*) = g^*$.
This is precisely the diagonal fixed-point equation solved above.
The unique non-trivial positive real fixed point is $g^* = \varphi$.

\textbf{Physical meaning.}
The renormalization group governs how a coupling constant changes with energy scale.
A fixed point is an energy scale at which the coupling stops running --- it is
self-consistent, in the sense that using the coupling to evolve itself brings
you back to the same value.
In the SM, the Higgs VEV $v$ corresponds to the scale at which the
Higgs self-coupling passes through zero.
In the SRRG, the MDL-minimizing fixed point at $g^* = \varphi$ identifies
the physical electroweak scale.
The physical meaning of this self-consistency is direct: the universe's
internal description (the polynomial applied to itself) must be consistent with
itself.
The coupling value at which this holds is $\varphi$, and the energy scale
associated with that coupling is 246~GeV --- the electroweak scale.
No parameter is adjusted to achieve this; it is forced by the polynomial's structure.
The theorem \lean{op9\_catal\_unconditional} (\CatAL, zero sorry, zero custom axioms)
establishes this identification without any SM input~\cite{SpivackSRRG,SpivackGTECompleteFramework}.

%─────────────────────────────────────────────────────────────────────────────
\section{From the Golden Ratio to the Higgs VEV}
\label{sec:vev}
%─────────────────────────────────────────────────────────────────────────────

\subsection{The derivation step by step}

The SRRG fixed point $g^* = \varphi$ sets the coupling at the electroweak scale.
To extract the Higgs VEV $v_{\rm PSC}$ from this coupling,
the SRRG scaling relation identifies:
\begin{equation}
  v_{\rm PSC} = 246.16~\text{GeV} \qquad (-0.024\%~\text{from PDG } 246.22~\text{GeV};\ \CatAL).
  \label{eq:vev}
\end{equation}
Lean cert: \lean{op9\_catal\_unconditional}~\cite{SpivackSRRG}.

No SM parameter enters this derivation.
The Fermi constant $G_F$ (from which the SM extracts $v$) is not used.
The coupling between the polynomial's arithmetic and the physical VEV is
the content of the \CatAL\ theorem.

\begin{tcolorbox}[colback=boxorange,colframe=orange!60!black,
  title={Result: Higgs VEV from SRRG (\CatAL)}]
  $v_{\rm PSC} = 246.16$~GeV\\[2pt]
  PDG~2024 (from $G_F$): $v = 246.22$~GeV\\[2pt]
  Deviation: $\mathbf{-0.024\%}$. Zero free parameters.\\[2pt]
  Lean: \lean{op9\_catal\_unconditional}~\cite{SpivackSRRG,SpivackGTECompleteFramework}
\end{tcolorbox}

\subsection{Why no fine-tuning}

The key point is that $v_{\rm PSC} = 246.16$~GeV is not the result of
cancelling two large numbers.
It is the unique value forced by the self-referential fixed-point equation
$p(g^*,g^*,g^*) = g^*$.

The 17-order-of-magnitude gap between $v$ and $M_{\rm Pl}$ is real.
But in the GTE framework, this gap does not arise from a cancellation;
it arises from the difference between the electroweak fixed point
($g^* = \varphi \approx 0.618$) and the Planck-scale structure
($M_{\rm Pl}/m_\tau = F_{21}^{10} \times |\mathbb{Z}_7|^{7/2}$, $0.040\%$
precision~\cite{SpivackGTECompleteFramework}).
Both are theorems from the same polynomial.
Their ratio is a theorem too.
No fine-tuning, no cancellation --- the hierarchy is a structural consequence
of the polynomial's arithmetic.

%─────────────────────────────────────────────────────────────────────────────
\section{From the VEV to the Higgs Mass}
\label{sec:mass}
%─────────────────────────────────────────────────────────────────────────────

\subsection{The Higgs quartic coupling}

In the SM, the Higgs potential is $V(\phi) = -\mu^2 |\phi|^2 + \lambda |\phi|^4$.
The Higgs VEV is $v = \sqrt{\mu^2/\lambda}$, and the Higgs mass is
$m_H = \sqrt{2\lambda}\, v$.
Both $\mu$ and $\lambda$ are free parameters in the SM.

In the GTE framework, the quartic coupling $\lambda$ is derived from the
SRRG golden-ratio fixed point with an information-processing-time (IPT)
correction~\cite{SpivackGTECompleteFramework}:
\begin{equation}
  \lambda = \frac{\varphi}{4\pi}
  \left(1 + \frac{\mathrm{IPT} - 1}{2c_H + 1}\right),
  \label{eq:lambda}
\end{equation}
where:
\begin{itemize}[itemsep=2pt]
  \item $\varphi = (\sqrt{5}-1)/2$ is the golden ratio (SRRG fixed point).
  \item $\mathrm{IPT} = 1 + \ln\varphi/(2\ln 2\pi)$ is the SRRG
        information-processing time.
  \item $c_H = 13$ is the Higgs color count: the number of distinct
        $\mathbb{Z}_7$ orbit sectors accessible to the Higgs doublet
        under the $F_{21}$ action.
  \item $2c_H + 1 = 27 = N_{\rm gen}^3 = 3^3$, certified as an algebraic
        identity by the theorem \lean{two\_cH\_plus\_one\_eq\_ngen\_cubed}
        (\CatAL)~\cite{SpivackGTECompleteFramework}.
\end{itemize}

\subsection{Worked calculation of the Higgs mass}

\textbf{Step 1.} The golden ratio:
\[
  \varphi = \frac{\sqrt{5}-1}{2} = 0.61803\ldots
\]

\textbf{Step 2.} The IPT correction:
\[
  \mathrm{IPT} = 1 + \frac{\ln\varphi}{2\ln(2\pi)}
  = 1 + \frac{\ln(0.61803)}{2\ln(6.2832)}
  = 1 + \frac{-0.48121}{3.6920}
  = 1 - 0.13035
  = 0.86965.
\]

\textbf{Step 3.} The denominator $2c_H + 1 = 2(13) + 1 = 27$.

\textbf{Step 4.} The coupling:
\begin{align*}
  \lambda &= \frac{0.61803}{4\pi}
  \left(1 + \frac{0.86965 - 1}{27}\right)
  = \frac{0.61803}{12.566} \times \left(1 - \frac{0.13035}{27}\right)\\
  &= 0.049185 \times (1 - 0.004828)
  = 0.049185 \times 0.995172
  = 0.048947.
\end{align*}

\textbf{Step 5.} The Higgs mass using $v_{\rm PSC} = 246.16$~GeV:
\begin{align*}
  m_H &= \sqrt{2\lambda}\, v_{\rm PSC}
  = \sqrt{2 \times 0.048947} \times 246.16~\text{GeV}\\
  &= \sqrt{0.097894} \times 246.16~\text{GeV}
  = 0.31288 \times 246.16~\text{GeV}
  = 125.25~\text{GeV}.
\end{align*}

\begin{tcolorbox}[colback=boxorange,colframe=orange!60!black,
  title={Result: Higgs Mass (\CatAD)}]
  $m_H^{\rm GTE} = 125.2499$~GeV\\[2pt]
  PDG~2024: $m_H = 125.20 \pm 0.11$~GeV\\[2pt]
  Deviation: $\mathbf{+0.45\sigma}$. Zero free parameters.\\[2pt]
  Key algebraic cert: \lean{two\_cH\_plus\_one\_eq\_ngen\_cubed} (\CatAL)~\cite{SpivackGTECompleteFramework}
\end{tcolorbox}

\subsection{The self-consistency requirement}

The Higgs mass calculation uses the self-consistent SRRG VEV $v_{\rm PSC}$,
not a tree-level estimate.
Without this self-consistency --- using an uncorrected $v$ instead of
$v_{\rm PSC}$ --- the prediction would be approximately $-2.08\sigma$
from PDG.
The $+0.45\sigma$ result follows from using the same fixed-point equation
that gives $v_{\rm PSC}$ to also correct the effective quartic coupling.
This is a single coherent system: the polynomial's fixed point determines
both the VEV and the mass together.

%─────────────────────────────────────────────────────────────────────────────
\section{Verification Against the Canonical Run}
\label{sec:verify}
%─────────────────────────────────────────────────────────────────────────────

The GTE SM derivation has been independently verified through a dual-path
computation.
The canonical run stores its results at
\texttt{papers/01\_SM/canonical\_run/dual\_path\_comparison.json}
in the GTE corpus.

That computation confirms:
\begin{itemize}[itemsep=2pt]
  \item Higgs VEV: $v_{\rm PSC} = 246.16$~GeV ($-0.024\%$ from PDG).
  \item Higgs mass: $m_H = 125.2499$~GeV ($+0.45\sigma$ from PDG).
  \item Both results are consistent across two independent derivation paths.
\end{itemize}

The canonical run uses no PDG Higgs data as input on either path.
Agreement between the two independent paths is a strong internal consistency check.

%─────────────────────────────────────────────────────────────────────────────
\section{Why No SUSY, No New Physics Below $\Lambda_{\rm GTE}$}
\label{sec:nosusy}
%─────────────────────────────────────────────────────────────────────────────

The GTE solution to the hierarchy problem makes sharp falsifiable predictions
about what \emph{not} to find.

\subsection{No supersymmetric particles}

The SRRG fixed-point mechanism resolves fine-tuning without introducing
any sparticles.
The electroweak scale is set by the polynomial's own arithmetic, not by
cancellation between SM particles and their superpartners.
If SUSY is not responsible for stabilizing the Higgs mass, there is no
reason to expect SUSY particles at any energy scale accessible to collider experiments.
GTE predicts no sparticles~\cite{SpivackGTECompleteFramework}.

\subsection{No new physics below $\Lambda_{\rm GTE}$}

The GTE framework has a single intrinsic scale beyond the SM:
\[
  \Lambda_{\rm GTE} = 7 M_{\rm kink} \approx 2.0~\text{GeV},
\]
the seven-kink full-winding threshold.
This is derived, not fitted ($\CatAD$; tree-level reading
$(8/7)m_\tau = 2.031$~GeV).
No new physics is predicted between the electroweak scale and $\Lambda_{\rm GTE}$
(i.e., between $\sim 100$~GeV and $\sim 2$~GeV, going downward in energy).
In particular: no axion, no Kaluza--Klein modes, no light dark-sector
resonances in the $1$--$100$~GeV window.

\subsection{The falsification profile}

\begin{tcolorbox}[colback=boxblue,colframe=blue!40!black,
  title={Falsifiable predictions from the SRRG hierarchy solution}]
\begin{itemize}[itemsep=2pt]
  \item No sparticles at any accessible collider energy (\CatAD).
  \item No axion (\CatAL, from strong CP theorem
    \lean{f21\_theta\_term\_vanishes}).
  \item Higgs mass $m_H = 125.2499$~GeV ($+0.45\sigma$):
    a future more precise Higgs mass measurement that excludes this value
    at $5\sigma$ would falsify the SRRG derivation.
  \item Higgs VEV $v = 246.16$~GeV:
    a future $G_F$ determination placing $v$ outside the GTE-predicted
    range at $5\sigma$ would falsify the derivation.
\end{itemize}
\end{tcolorbox}

%─────────────────────────────────────────────────────────────────────────────
\section{The Gravitational Hierarchy: Also a Theorem}
\label{sec:gravity}
%─────────────────────────────────────────────────────────────────────────────

The original hierarchy problem concerned the Higgs mass.
But the deepest hierarchy in physics is gravitational: gravity is
$10^{32}$ times weaker than the weak force.
In the SM, Newton's constant $G_N$ is a free parameter.
Why should gravity be so extraordinarily weak?

In the GTE framework, this ratio too is a theorem.
The Planck-to-tau-mass ratio is:
\[
  \frac{M_{\rm Pl}}{m_\tau} = F_{21}^{10} \times |\mathbb{Z}_7|^{7/2},
\]
where $F_{21} = \mathbb{Z}_7 \rtimes \mathbb{Z}_3$ is the order-21
Frobenius group whose arithmetic generates the GTE polynomial,
and $|\mathbb{Z}_7| = 7$.
This gives $M_{\rm Pl}/m_\tau \approx 6.87 \times 10^{18}$, matching
the measured ratio ($M_{\rm Pl}/m_\tau = 6.88 \times 10^{18}$) at
$0.040\%$ precision (\CatAD)~\cite{SpivackGTECompleteFramework}.

The gravitational hierarchy is not a fine-tuning.
It is a consequence of group theory.

\begin{tcolorbox}[colback=boxgold,colframe=orange!50!black,
  title={The hierarchy problem, dissolved}]
The hierarchy problem asks why the electroweak scale ($\sim 100$~GeV) is
so much smaller than the Planck scale ($\sim 10^{19}$~GeV).

In the GTE framework, both scales are derived from the same 19-bit
polynomial $p(L,C,R) = C + R - CR - LCR \pmod{7}$:
\begin{itemize}[itemsep=2pt]
  \item \textbf{Electroweak scale:} SRRG fixed point $p(x,x,x) = x$
        gives $x = \varphi$ (golden ratio), setting $v = 246.16$~GeV
        (\CatAL).
  \item \textbf{Planck scale:} $M_{\rm Pl}/m_\tau = F_{21}^{10}
        \times 7^{7/2}$ from $F_{21}$ orbit combinatorics (\CatAD).
  \item \textbf{Their ratio:} a group-theoretic theorem, not a
        fine-tuning.
\end{itemize}
No cancellation. No sparticles. No anthropic selection.
The hierarchy is structural.
\end{tcolorbox}

%─────────────────────────────────────────────────────────────────────────────
\section*{Further Reading}
%─────────────────────────────────────────────────────────────────────────────
\begin{tcolorbox}[colback=orange!8,colframe=orange!50!black,
  title=\textbf{Primary Sources}]
The results in this tutorial are derived in the following papers,
all available open-access on Zenodo and at
\href{https://novaspivack.com/research}{novaspivack.com/research}:
\begin{itemize}[itemsep=2pt]
  \item \textbf{P48 --- The Complete GTE Framework} (main synthesis monograph):\\
        \href{https://doi.org/10.5281/zenodo.20560550}{doi.org/10.5281/zenodo.20560550}
  \item \textbf{P47 --- Cosmological Predictions of the GTE/$\Phi_{\rm MDL}$ Framework}
        (SRRG and Higgs sector derivations):\\
        \href{https://doi.org/10.5281/zenodo.20465803}{doi.org/10.5281/zenodo.20465803}
  \item \textbf{P01 --- Standard Model Parameters from the UGP} (full SM derivation):\\
        \href{https://doi.org/10.5281/zenodo.20168787}{doi.org/10.5281/zenodo.20168787}
\end{itemize}
Lean~4 machine proofs: \href{https://github.com/novaspivack/ugp-lean}{github.com/novaspivack/ugp-lean}
\end{tcolorbox}

\bibliographystyle{unsrtnat}
\bibliography{../../papers/bib/Spivack_Papers_Bibliography}

\end{document}
